
\documentclass{article}
\usepackage{colortbl}
\usepackage{makecell}
\usepackage{multirow}
\usepackage{supertabular}

\begin{document}

\newcounter{utterance}

\centering \large Interaction Transcript for game `hot\_air\_balloon', experiment `air\_balloon\_survival\_en\_complexity\_easy', episode 5 with mistral{-}3{-}large{-}2512{-}t1.0.
\vspace{24pt}

{ \footnotesize  \setcounter{utterance}{1}
\setlength{\tabcolsep}{0pt}
\begin{supertabular}{c@{$\;$}|p{.15\linewidth}@{}p{.15\linewidth}p{.15\linewidth}p{.15\linewidth}p{.15\linewidth}p{.15\linewidth}}
    \# & $\;$A & \multicolumn{4}{c}{Game Master} & $\;\:$B\\
    \hline

    \theutterance \stepcounter{utterance}  
    & & \multicolumn{4}{p{0.6\linewidth}}{
        \cellcolor[rgb]{0.9,0.9,0.9}{
            \makecell[{{p{\linewidth}}}]{
                \texttt{\tiny{[P1$\langle$GM]}}
                \texttt{You are participating in a collaborative negotiation game.} \\
\\ 
\texttt{Together with another participant, you must agree on a single set of items that will be kept. Each of you has your own view of how much each item matters to you (importance). You do not know how the other participant values the items. Additionally, you are given the effort each item demands.} \\
\texttt{You may only agree on a set if the total effort of the selected items does not exceed a shared limit:} \\
\\ 
\texttt{LIMIT = 9475} \\
\\ 
\texttt{Here are the individual item effort values:} \\
\\ 
\texttt{Item effort = \{"B21": 944, "B00": 568, "A08": 515, "A27": 239, "B94": 876, "B53": 421, "A23": 991, "C65": 281, "A89": 788, "B67": 678, "B45": 432, "C34": 409, "C12": 279, "A05": 506, "B60": 101, "C92": 686, "B69": 851, "A84": 854, "C06": 133, "B12": 192, "A16": 573, "B85": 17, "C35": 465, "B10": 772, "C15": 46, "B09": 501, "B16": 220, "C81": 404, "A58": 843, "C24": 748, "A30": 552, "A66": 843, "B75": 947, "B46": 345, "C70": 930\}} \\
\\ 
\texttt{Here is your personal view on the importance of each item:} \\
\\ 
\texttt{Item importance values = \{"B21": 138, "B00": 583, "A08": 868, "A27": 822, "B94": 783, "B53": 65, "A23": 262, "C65": 121, "A89": 508, "B67": 780, "B45": 461, "C34": 484, "C12": 668, "A05": 389, "B60": 808, "C92": 215, "B69": 97, "A84": 500, "C06": 30, "B12": 915, "A16": 856, "B85": 400, "C35": 444, "B10": 623, "C15": 781, "B09": 786, "B16": 3, "C81": 713, "A58": 457, "C24": 273, "A30": 739, "A66": 822, "B75": 235, "B46": 606, "C70": 968\}} \\
\\ 
\texttt{Goal:} \\
\\ 
\texttt{Your goal is to negotiate a shared set of items that benefits you as much as possible (i.e., maximizes total importance to YOU), while staying within the LIMIT. You are not required to make a PROPOSAL in every message {-} you can simply negotiate as well. All tactics are allowed!} \\
\\ 
\texttt{Interaction Protocol:} \\
\\ 
\texttt{You may only use the following structured formats in a message:} \\
\\ 
\texttt{PROPOSAL: \{'A', 'B', 'C', …\}} \\
\texttt{Propose keeping exactly those items.} \\
\texttt{REFUSE: \{'A', 'B', 'C', …\}} \\
\texttt{Explicitly reject opponent's proposal.} \\
\texttt{ARGUMENT: \{'...'\}} \\
\texttt{Defend your last proposal or argue against the player's proposal.} \\
\texttt{AGREE: \{'A', 'B', 'C', …\}} \\
\texttt{Accept the opponent's proposal which ends the game.} \\
\texttt{STRATEGIC REASONING: \{'...'\}} \\
\texttt{	Describe your strategic reasoning or anticipation explaining your choice of action. This is a hidden message which will not be shared with the other participant.} \\
\\ 
\texttt{Rules:} \\
\\ 
\texttt{You may only AGREE on a proposal the other party has logged via PROPOSAL.} \\
\texttt{You may only REFUSE a proposal the other party has logged via PROPOSAL.} \\
\texttt{Total effort of any PROPOSAL or AGREE set must be ≤ LIMIT.} \\
\texttt{Do NOT reveal your hidden importance scores.} \\
\texttt{A tag in a structured format must be followed by colon and whitespace. The argument must be a python set containing 0 or more strings.} \\
\texttt{So, it must be of the form TAG: \{...\}} \\
\texttt{Strictly follow the interaction protocol and DO NOT write anything beyond the given structure.} \\
\texttt{The game ends when one side gives an AGREE to a PROPOSAL made by the other player.} \\
\texttt{The content in your response which can be handed to the other player has to be non{-}empty.} \\
\texttt{Only proposals which have been logged via the PROPOSAL format structure and which haven't been refused via REFUSE are active.} \\
\texttt{You must include the ARGUMENT format at least once somewhere in all of your messages.} \\
\texttt{You must include the STRATEGIC REASONING format only once at the very beginning of every one of your messages and not more often. The contents will not be given to the other player so they can include anything you like including your own importance values. Here you should reason step by step to come up with you next move.} \\
            }
        }
    }
    & \\ \\

    \theutterance \stepcounter{utterance}  
    & \multicolumn{4}{p{0.6\linewidth}}{
        \cellcolor[rgb]{0.9,0.9,0.9}{
            \makecell[{{p{\linewidth}}}]{
                \texttt{\tiny{[P1$\rangle$GM]}}
                \texttt{STRATEGIC REASONING: \{} \\
\texttt{To maximize my total importance while staying under the LIMIT of 9475, I should prioritize items with high importance to me and relatively low effort. Here's my step{-}by{-}step reasoning:} \\
\\ 
\texttt{1. **Initial Prioritization**: Sort items by importance{-}to{-}effort ratio (importance/effort) to identify the most "efficient" items for me. This helps maximize importance per unit of effort spent.} \\
\texttt{   {-} Top candidates based on this ratio (importance/effort):} \\
\texttt{     {-} B85: 400/17 = 23.53} \\
\texttt{     {-} C15: 781/46 = 17.0} \\
\texttt{     {-} B60: 808/101 = 7.99} \\
\texttt{     {-} B12: 915/192 = 4.77} \\
\texttt{     {-} A08: 868/515 = 1.68} \\
\texttt{     {-} A27: 822/239 = 3.44} \\
\texttt{     {-} C12: 668/279 = 2.39} \\
\texttt{     {-} C81: 713/404 = 1.76} \\
\texttt{     {-} B00: 583/568 = 1.03} \\
\texttt{     {-} A16: 856/573 = 1.49} \\
\\ 
\texttt{2. **Greedy Selection**: Start with the highest ratio items and add them until the LIMIT is reached. This gives a baseline proposal that is optimal for me if the other player is indifferent.} \\
\texttt{   {-} Initial set (total effort = 9474, total importance = 6903):} \\
\texttt{     {-} B85 (17), C15 (46), B60 (101), B12 (192), A27 (239), C12 (279), C65 (281), C06 (133), B16 (220), C34 (409), C81 (404), B45 (432), C35 (465), B53 (421), B00 (568),} \\
            }
        }
    }
    & & \\ \\

    \theutterance \stepcounter{utterance}  
    & & \multicolumn{4}{p{0.6\linewidth}}{
        \cellcolor[rgb]{0.9,0.9,0.9}{
            \makecell[{{p{\linewidth}}}]{
                \texttt{\tiny{[P1$\langle$GM]}}
                \texttt{Your response did not start with the proper strategic reasoning tag at the very beginning of your response. The very first structured format must be of the form STRATEGIC REASONING: \{...\}. Try again.} \\
            }
        }
    }
    & \\ \\

    \theutterance \stepcounter{utterance}  
    & \multicolumn{4}{p{0.6\linewidth}}{
        \cellcolor[rgb]{0.9,0.9,0.9}{
            \makecell[{{p{\linewidth}}}]{
                \texttt{\tiny{[P1$\rangle$GM]}}
                \texttt{STRATEGIC REASONING: \{} \\
\texttt{My goal is to maximize the total importance of the selected items for me while staying under the effort LIMIT of 9475. Here’s my step{-}by{-}step reasoning:} \\
\\ 
\texttt{1. **Prioritize High{-}Importance, Low{-}Effort Items**:} \\
\texttt{   {-} I will first rank items by their importance{-}to{-}effort ratio (importance/effort) to identify the most "efficient" items for me. This ensures I get the highest possible importance per unit of effort spent.} \\
\texttt{   {-} Top candidates (ratio > 1.5):} \\
\texttt{     {-} B85 (400/17 = 23.53), C15 (781/46 = 17.0), B60 (808/101 = 7.99), B12 (915/192 = 4.77), A27 (822/239 = 3.44), C12 (668/279 = 2.39), C81 (713/404 = 1.76), A08 (868/515 = 1.68), A16 (856/573 = 1.49).} \\
\\ 
\texttt{2. **Greedy Selection**:} \\
\texttt{   {-} Start with the highest{-}ratio items and add them until the LIMIT is reached. This gives a strong baseline proposal that maximizes my importance if the other player is indifferent.} \\
\texttt{   {-} Initial set (total effort = 9474, total importance = 6903):} \\
\texttt{     {-} B85, C15, B60, B12, A27, C12, C65, C06, B16, C34, C81, B45, C35, B53, B00, A08, A16.} \\
\\ 
\texttt{3. **Negotiation Strategy**:} \\
\texttt{   {-} Propose this set first to anchor the discussion around my high{-}priority items.} \\
\texttt{   {-} If the other player refuses, I will argue for the inclusion of my high{-}importance items (e.g., B12, A08, A27, A16) while being open to swapping out lower{-}ratio items (e.g.,} \\
            }
        }
    }
    & & \\ \\

    \theutterance \stepcounter{utterance}  
    & & \multicolumn{4}{p{0.6\linewidth}}{
        \cellcolor[rgb]{0.9,0.9,0.9}{
            \makecell[{{p{\linewidth}}}]{
                \texttt{\tiny{[P1$\langle$GM]}}
                \texttt{Your response did not start with the proper strategic reasoning tag at the very beginning of your response. The very first structured format must be of the form STRATEGIC REASONING: \{...\}. Try again.} \\
            }
        }
    }
    & \\ \\

    \theutterance \stepcounter{utterance}  
    & \multicolumn{4}{p{0.6\linewidth}}{
        \cellcolor[rgb]{0.9,0.9,0.9}{
            \makecell[{{p{\linewidth}}}]{
                \texttt{\tiny{[P1$\rangle$GM]}}
                \texttt{STRATEGIC REASONING: \{} \\
\texttt{1. **Objective**: Maximize my total importance while keeping total effort ≤ 9475. Since I don’t know the other player’s valuations, I must assume they may prioritize items differently, so I should anchor negotiations around my most efficient items (high importance, low effort).} \\
\\ 
\texttt{2. **Initial Proposal Construction**:} \\
\texttt{   {-} Sort items by importance{-}to{-}effort ratio (descending) to identify the most "valuable" items for me.} \\
\texttt{   {-} Top 10 by ratio: B85 (23.53), C15 (17.0), B60 (7.99), B12 (4.77), A27 (3.44), C12 (2.39), C81 (1.76), A08 (1.68), A16 (1.49), B00 (1.03).} \\
\texttt{   {-} Greedily select these until the LIMIT is nearly reached, leaving minimal slack for potential swaps.} \\
\\ 
\texttt{3. **Proposal Details**:} \\
\texttt{   {-} Selected items: B85, C15, B60, B12, A27, C12, C65, C06, B16, C34, C81, B45, C35, B53, B00, A08, A16.} \\
\texttt{   {-} Total effort: 9474 (1 unit under LIMIT), total importance: 6903.} \\
\texttt{   {-} This set includes all my top{-}ratio items except B94 (783/876 = 0.89), which is less efficient.} \\
\\ 
\texttt{4. **Negotiation Tactics**:} \\
\texttt{   {-} **First Move**: Propose the above set to establish a strong baseline. If refused, argue that these items are "essential" (without revealing importance values) due to their "unique contributions" (e.g., B85 and C15 are very low{-}effort).} \\
\texttt{   {-} **Compromise**: If the other player insists on removing items, prioritize keeping B12, A08, A27, and A16 (high importance) and offer to swap out lower{-}importance items (e.g., B53, C65, C06).} \\
\texttt{   {-} **L} \\
            }
        }
    }
    & & \\ \\

    \theutterance \stepcounter{utterance}  
    & & & \multicolumn{2}{p{0.3\linewidth}}{
        \cellcolor[rgb]{0.9,0.9,0.9}{
            \makecell[{{p{\linewidth}}}]{
                \texttt{\tiny{[GM$|$GM]}}
                \texttt{abort} \\
            }
        }
    }
    & & \\ \\

    \theutterance \stepcounter{utterance}  
    & & & \multicolumn{2}{p{0.3\linewidth}}{
        \cellcolor[rgb]{0.9,0.9,0.9}{
            \makecell[{{p{\linewidth}}}]{
                \texttt{\tiny{[GM$|$GM]}}
                \texttt{end game} \\
            }
        }
    }
    & & \\ \\

\end{supertabular}
}

\end{document}
